\chapter*{译者说明}

这份文稿是丁鹏《A First Course in Causal Inference》的中文译本工程。第一阶段的目标不是一次性完成全书翻译，而是先建立一个稳定、可编译、可持续扩展的中文教材骨架。

本译本遵循“忠实翻译 + 注解”的原则：数学定义、命题、公式、例子、图表、引用和原书标签尽量保持原结构；中文正文则尽量使用自然、清楚、适合学习的表达。专业术语在首次出现时采用“中文术语（English term）”的形式，重要缩写保留英文，例如 Fisher 随机化检验（Fisher randomization test, FRT）。\footnote{关于为什么采用英文括注，见附录 \cref{sec:qa-english-term-parentheses}。}

读者在阅读过程中提出的问题会记录到问答附录，并在正文的自然位置通过脚注回链。这样，这本中文教材会随着阅读对话逐渐生长，而不是停留在一次性的静态翻译。
